\section{AGI 的定义与时间线}

\subsection{什么是 AGI}

Demis Hassabis 将 AGI（Artificial General Intelligence）定义为\textbf{能够模拟人类全部认知能力的人工智能系统}。与当前的``窄 AI''不同，AGI 需要具备跨领域的通用推理能力——不仅仅是在某一项任务上表现出色，而是能够像人类大脑一样，将来自不同领域的知识编织成连贯的整体。

\begin{importantbox}{AGI 的核心定义}
AGI = 能够模拟人类全部认知能力（full human cognitive capabilities）的通用人工智能。关键特征包括：
\begin{itemize}
    \item 跨领域推理与知识整合
    \item 持续学习与记忆更新
    \item 长期规划与多步骤决策
    \item 一致性的通用推理能力
\end{itemize}
这不是某个单一 benchmark 的突破，而是认知能力的\textbf{全面覆盖}。
\end{importantbox}

\subsection{时间线：五年之内}

Hassabis 明确表示，AGI 在\textbf{未来五年内实现的可能性很大}（``very good chance of it being within the next 5 years''）。值得注意的是，这一判断与 DeepMind 联合创始人 Shane Legg 在 2010 年的预测一脉相承——当时 Legg 预测 AGI 将在 20 年内到来，而现在距离那个预测恰好过去了约 16 年。

驱动 AGI 到来的两大引擎是：
\begin{enumerate}
    \item \textbf{算力规模化（Compute Scaling）}：更大的计算集群、更高效的硬件
    \item \textbf{算法进步（Algorithmic Progress）}：新的架构设计、训练方法与优化策略
\end{enumerate}

\begin{knowledgebox}{Shane Legg 的远见}
Shane Legg 是 DeepMind 的联合创始人之一，早在 2010 年就公开预测 AGI 将在 20 年内实现。当时这被视为极端乐观的预测，但如今 Hassabis 指出这条时间线正在兑现。从最初的三人创业（Hassabis、Legg 与 Mustafa Suleyman）到今天的 Google DeepMind，这种对 AGI 的长期信念一直是团队的核心驱动力。
\end{knowledgebox}

\subsection{量化 AGI 的影响}

Hassabis 给出了一个令人震撼的量化表述：

\begin{importantbox}{``十倍工业革命、十倍速度''}
\textit{``我有时将 AGI 的到来量化为十倍的工业革命，以十倍的速度发生。''（``I sometimes quantify the coming of AGI as 10 times the industrial revolution at 10 times the speed.''）}

工业革命用了一百多年才让人类消化其副作用（童工、城市化阵痛、阶级冲突等），而 AGI 带来的变革将在十年内压缩完成。这意味着人类社会几乎没有缓冲时间来适应。
\end{importantbox}

工业革命前，儿童死亡率高达 40\%；工业革命后的科技进步让人类寿命翻倍、生活水平大幅提升。但这一过程历时逾百年。AGI 的承诺是将同等规模（甚至更大）的变革压缩进一代人的生命周期——这既是巨大的机遇，也是前所未有的挑战。

\subsection{本章小结}

Hassabis 将 AGI 定义为模拟人类全部认知能力的通用智能，预期五年内实现。他用``十倍工业革命、十倍速度''来量化其影响规模，强调人类社会必须以前所未有的紧迫感来准备这一变革。

%% ==================== Section 2 ====================
